\subsection{Лабораторная работа № 1}

\subsubsection{Задача}

В ходе лабораторной работы необходимо:

\begin{enumerate}
    \item создать представление, содержащее результат группировки данных с
    двумя агрегатными функциями \mintinline{sql}{COUNT};
    \item составить простой запрос к созданному представлению с условием
    \mintinline{sql}{WHERE};
    \item составить запрос, в котором представление используется совместно с
    операциями соединения \mintinline{sql}{JOIN}.
\end{enumerate}

\subsubsection{Создание представления}

Необходимо создать представление, которое для каждого начертания хранит число
поддерживаемых форматов и количество градаций насыщенности.

Одно начертание может храниться в таблице \mintinline{text}{Typeface} в
нескольких форматах и градациях насыщенности. Чтобы не выполнять одинаковый
подсчёт в каждом запросе, создаётся отдельное представление.

В нём для каждого начертания хранятся название, идентификатор гарнитуры, число
форматов и число градаций насыщенности. Определение представления приведено в
листинге~\ref{lst:create-view}.

В формулах используются обозначения: $T$~--- таблица
\mintinline{text}{Typeface}, $F$~--- таблица
\mintinline{text}{Format}, $W$~--- таблица
\mintinline{text}{Weight}, $FF$~--- таблица
\mintinline{text}{Font_family}, $C$~--- таблица
\mintinline{text}{Category}, $V$~--- созданное представление. Символ
$\bowtie_{k}$ обозначает соединение по полю $k$, $\sigma$~---
выборку, $\pi$~--- проекцию, $\gamma$~--- группировку с агрегированием.

\codeListing[sql]{../practice-01/src/create.sql}{Создание представления}{\label{lst:create-view}}

Агрегатная функция \mintinline{sql}{COUNT(DISTINCT ...)} подсчитывает
количество различных значений. Поэтому повторяющиеся форматы и градации
насыщенности учитываются только один раз. Для получения фактических форматов
и градаций насыщенности таблица \mintinline{text}{Typeface} соединяется с
таблицами \mintinline{text}{Format} и \mintinline{text}{Weight}.

Создание представления описывается формулой

\[
V = \gamma_{\substack{id\_font\_family,name;\\
\operatorname{count\_distinct}(id\_format) \rightarrow format\_count,\\
\operatorname{count\_distinct}(id\_weight) \rightarrow weight\_count}}
\left(
    \left(T \bowtie_{id\_format} F\right)
    \bowtie_{id\_weight} W
\right).
\]

\subsubsection{Проверка представления}

\textit{Текст запроса:} Вывести первые десять начертаний из созданного представления вместе с числом
поддерживаемых форматов и градаций насыщенности.

Запрос приведён в листинге~\ref{lst:query-01}.

\codeListing[sql]{../practice-01/src/query-01.sql}{Проверка содержимого представления}{\label{lst:query-01}}

Реляционная формула имеет вид

\[
R_1 = \pi_{\substack{id\_font\_family,typeface\_name,\\
format\_count,weight\_count}}(V).
\]

План и результат выполнения показаны на
рисунках~\ref{fig:query-01-plan} и~\ref{fig:query-01-result}.

\img[0.85]{img/pr-01-query-01-plan.png}{План выполнения запроса 1}{\label{fig:query-01-plan}}

PostgreSQL читает таблицы \mintinline{text}{Typeface},
\mintinline{text}{Format} и \mintinline{text}{Weight}, затем соединяет их по
полям \mintinline{sql}{id_format} и \mintinline{sql}{id_weight}. Полученные
строки сортируются по полям группировки, после чего узел
\mintinline{text}{GroupAggregate} рассчитывает количество различных форматов
и градаций насыщенности каждого начертания.

\img[0.8]{img/pr-01-query-01-result.png}{Результат выполнения запроса 1}{\label{fig:query-01-result}}

\clearpage
\subsubsection{Запрос с условием}

\textit{Текст запроса:} Найти число начертаний, у которых количество поддерживаемых форматов больше
$A$.

В запросе значение $A$ принято равным 2. Сначала условие
\mintinline{sql}{format_count > 2} оставляет подходящие начертания, затем
\mintinline{sql}{COUNT} подсчитывает их. Запрос приведён в
листинге~\ref{lst:query-02}.

\codeListing[sql]{../practice-01/src/query-02.sql}{Подсчёт начертаний с числом форматов больше
$A$}{\label{lst:query-02}}

Запрос описывается формулой

\[
R_2 = \gamma_{\operatorname{count}(id\_font\_family)
\rightarrow typeface\_count}
\left(\sigma_{format\_count>A}(V)\right).
\]

План и результат выполнения показаны на
рисунках~\ref{fig:query-02-plan} и~\ref{fig:query-02-result}.

\img[0.8]{img/pr-01-query-02-plan.png}{План выполнения запроса 2}{\label{fig:query-02-plan}}

Сначала PostgreSQL соединяет таблицу начертаний с таблицами форматов и
насыщенностей, после чего вычисляет строки представления с помощью сортировки
и \mintinline{text}{GroupAggregate}. Во время агрегирования отбрасываются
начертания, у которых число форматов не больше $A$. Затем узел
\mintinline{text}{Aggregate} подсчитывает оставшиеся строки.

\img[0.7]{img/pr-01-query-02-result.png}{Результат выполнения запроса 2}{\label{fig:query-02-result}}

Условию соответствует 39 начертаний.

\subsubsection{Охват форматов по категориям}

\textit{Текст запроса:} Найти пять категорий с наибольшим суммарным охватом поддерживаемых форматов.

Представление содержит идентификатор гарнитуры, поэтому его можно соединить с
таблицами \mintinline{text}{Font_family} и \mintinline{text}{Category}. После
соединения значения \mintinline{sql}{format_count} суммируются отдельно для
каждой категории. Запрос приведён в листинге~\ref{lst:query-03}.

\codeListing[sql]{../practice-01/src/query-03.sql}{Первые пять категорий по охвату
поддерживаемых форматов}{\label{lst:query-03}}

Агрегирование по категориям и выбор первых пяти строк описываются формулой

\begingroup
\small
\[
\begin{aligned}
R_3 ={}& \operatorname{top}_{5}^{format\_count\downarrow,\,name\uparrow}
\biggl(\\
&
    \gamma_{\substack{id\_category,name;\\
    \operatorname{sum}(format\_count) \rightarrow format\_count}}
    \left(
        \left(V \bowtie_{id\_font\_family} FF\right)
        \bowtie_{id\_category} C
    \right)
\biggr).
\end{aligned}
\]
\endgroup

Представление уже содержит число форматов каждого начертания, поэтому повторно
считать их в запросе не требуется. План и результат выполнения показаны на
рисунках~\ref{fig:query-03-plan} и~\ref{fig:query-03-result}.

\img[0.65]{img/pr-01-query-03-plan.png}{План выполнения запроса 3}{\label{fig:query-03-plan}}

PostgreSQL вычисляет представление, соединяя начертания с форматами и
насыщенностями. Полученные строки дополнительно соединяются с гарнитурами и
категориями. Затем значения \mintinline{sql}{format_count} суммируются по
категориям. Завершающие сортировка и ограничение оставляют пять категорий с
наибольшим охватом.

\img[0.75]{img/pr-01-query-03-result.png}{Результат выполнения запроса 3}{\label{fig:query-03-result}}

После группировки категории сортируются по убыванию охвата. Оператор
\mintinline{sql}{LIMIT 5} оставляет первые пять строк. Наибольший охват в
полученном результате имеет категория \mintinline{text}{Neo-grotesque}.
